\documentclass[reqno]{amsart} 
\usepackage{amssymb,latexsym,amsmath,amscd,graphicx,setspace,amsthm,verbatim,enumitem}
\usepackage[margin = 3 cm]{geometry}

\def\upint{\mathchoice%
    {\mkern13mu\overline{\vphantom{\intop}\mkern7mu}\mkern-20mu}%
    {\mkern7mu\overline{\vphantom{\intop}\mkern7mu}\mkern-14mu}%
    {\mkern7mu\overline{\vphantom{\intop}\mkern7mu}\mkern-14mu}%
    {\mkern7mu\overline{\vphantom{\intop}\mkern7mu}\mkern-14mu}%
  \int}
\def\lowint{\mkern3mu\underline{\vphantom{\intop}\mkern7mu}\mkern-10mu\int}


\theoremstyle{plain}
\newtheorem{theorem}{Theorem}[section]
\newtheorem{proposition}{Proposition}
\newtheorem{corollary}{Corollary}
\newtheorem{lemma}{Lemma}
\newtheorem{conjecture}{Conjecture}
\newtheorem{question}{Question}
\newtheorem{problem}{Problem}
      
\theoremstyle{definition}
\newtheorem{definition}{Definition}

\newenvironment{solution1}{\paragraph{\emph{Solution $1$}.}}{\hfill$\square$}
\newenvironment{solution2}{\paragraph{\emph{Solution $2$}.}}{\hfill$\square$}
\newenvironment{solution3}{\paragraph{\emph{Solution $3$}.}}{\hfill$\square$}

\begin{document} 

\title[Homework 9]{Homework 9}

\date{\today} 
\maketitle 

Recall from Homework $5$:  if $X$ is a set and $\mathcal{S} \subseteq \mathcal{P}(X)$, then the intersection of all $\sigma$-algebras on $X$ containing $\mathcal{S}$ is a $\sigma$-algebra on $X$ containing $\mathcal{S}$.  In fact, it is the smallest one containing $\mathcal{S}$.  We shall denote this $\sigma$-algebra by $\langle \mathcal{S} \rangle$.  With this notation, one has $\mathcal{B}(\mathbb{R}^{n}) = \langle \tau \rangle$, where $\tau$ is the collection of all open subsets of $\mathbb{R}^{n}$.


\begin{problem}
Let $(Y,\mathcal{A})$ be a measurable space and let $f: X \rightarrow Y$ be a function.  Assume that $\mathcal{A} = \langle \mathcal{S} \rangle$ for some $\mathcal{S} \subseteq \mathcal{P}(Y)$.  Show that
$$f^{*}\mathcal{A}= \langle \{f^{-1}(S) : S \in \mathcal{S} \}\rangle. $$
(The pull-back $f^{*}\mathcal{A}$ was defined in Problem $4$ of Homework $6$.)
\end{problem}
\begin{solution1}

\end{solution1}

\begin{problem}
Show that if $f:\mathbb{R}^{m} \rightarrow \mathbb{R}^{n}$ is a continuous function, then given any $B \in \mathcal{B}(\mathbb{R}^{n})$, one has $f^{-1}(B) \in \mathcal{B}(\mathbb{R}^{m})$.
\end{problem}
\begin{solution1}

\end{solution1}



\begin{problem}
Consider the Vitali set $V \subseteq [-1,1]$ and let $E = V \times \{0\} \subseteq \mathbb{R}^{2}$.  Does one have $E \in \mathcal{L}(\mathbb{R}^{2})$?  Explain carefully.
\end{problem}
\begin{solution1}

\end{solution1}


\begin{problem}
Consider the function $f: \mathbb{R} \rightarrow \mathbb{R}^{2}$ given by $x \mapsto f(x) = (x,0)$.  
\begin{enumerate}
\item Show that $f$ is a continuous function.
\item What is $f^{-1}(E)$, where $E$ was defined in the previous problem?
\item Does one have $E \in \mathcal{B}(\mathbb{R}^{2})$?  Explain carefully.
\item We know that $\mathcal{B}(\mathbb{R}^{2}) \subseteq \mathcal{L}(\mathbb{R}^{2})$, but do we have an equality?  Explain carefully.
\end{enumerate}
\end{problem}
\begin{solution1}

\end{solution1}


\begin{problem}
Is the measure space $(\mathbb{R}^{2},\mathcal{B}(\mathbb{R}^{2}),\lambda)$, where $\lambda$ is the Lebesgue measure restricted to the Borel sets $\mathcal{B}(\mathbb{R}^{2})$, a complete measure space??  Explain carefully.
\end{problem}
\begin{solution1}

\end{solution1}


The same statements as above are true for $\mathbb{R}$ instead of $\mathbb{R}^{2}$, but the proofs are a bit more complicated (at least the ones I know).  If you are interested you can go over the following {\bfseries optional} problems after reading about the Cantor function $\Lambda$ on page $58$ of your book.  (These problems below will {\bfseries not} appear on any exams.)

\noindent\rule{\linewidth}{0.4pt}

The goal is to come up with a set $A \in \mathcal{L}(\mathbb{R}) \smallsetminus \mathcal{B}(\mathbb{R})$.  This will be accomplished in several steps which I have broken down into problems with hints.  Throughout this problem set, given $S \subseteq \mathbb{R}$, we let 
$$S - S :=\{s_{1} - s_{2} : s_{1}, s_{2} \in S \}. $$

\begin{problem}
Show that given any $A \in \mathcal{L}(\mathbb{R})$ such that $\infty > \lambda(A) > 0$, there exists $\delta > 0$ such that
$$(-\delta,\delta) \subseteq A - A. $$
\end{problem}
Hint:  The idea relies on the observation that if $(t + A) \cap A \neq \varnothing$, then $t \in A - A$.  So we have to show that if $|t|$ is small, then $(t+A) \cap A \neq \varnothing$.  For that purpose, give yourself an $\varepsilon$ satisfying $0 < \varepsilon < 1$.  Argue first that there exists an open set $G$ of $\mathbb{R}$ such that $A \subseteq G$ and for which 
$$\lambda(G) < \lambda(A)(1+\varepsilon). $$
Then, write $G$ as a countable disjoint union of open intervals $I_{n}$.  Let $A_{n} = A \cap I_{n}$.  Show that the last displayed equation implies that there exists $k \in \mathbb{N}$ such that
$$\lambda(I_{k}) < (1+\varepsilon)\lambda(A_{k}). $$
For that $k$, explain why $I_{k}$ is necessarily a finite open interval, and then write $I_{k} = (a_{k},b_{k})$ for some real numbers $a_{k}, b_{k}$.  Explain why 
$$\lambda(I_{k} \cup (I_{k} + t)) \leq \lambda(I_{k}) + |t|. $$
At last, show that if $A_{k} \cap (t + A_{k}) = \varnothing$, then
$$2\lambda(I_{k})< (1+\varepsilon)(\lambda(I_{k}) + |t|). $$
Conclude the proof of Problem 6.

\begin{problem}
Show that if $A \in \mathcal{L}(\mathbb{R})$ is such that $\infty > \lambda(A) > 0$, then there exists $E \subseteq A$ such that $E \notin \mathcal{L}(\mathbb{R})$.
\end{problem}
Hint:  This is a refinement of our proof in class that there exist sets that are not Lebesgue measurable.  Define the relation $\sim$ on $\mathbb{R}$ via $x \sim y$ if $x - y \in \mathbb{Q}$.  Show that this relation is an equivalence relation.  Let $V$ be a complete set of representatives for this equivalence relation.  Show that
$$\mathbb{R} = \bigcup_{r \in \mathbb{Q}} (r+V),$$
and that this union is disjoint.  For $r \in \mathbb{Q}$, let $A_{r} = A \cap (r+V)$.  If $A_{r} \notin \mathcal{L}(\mathbb{R})$ for some $r \in \mathbb{Q}$, explain why we are done.  If $A_{r} \in \mathcal{L}(\mathbb{R})$ for all $r \in \mathbb{Q}$, then explain why
$$A_{r} - A_{r} \subseteq V - V \text{ and } (V - V) \cap \mathbb{Q}^{\times} = \varnothing. $$
Using Problem 6, show then that $\lambda(A) = 0$, a contradiction...

\begin{problem}
Consider the function $f:[0,1] \rightarrow [0,2]$ defined via
$$x \mapsto f(x) = \Lambda(x) + x, $$
where $\Lambda$ is the Cantor function.  Show that $f$ is a continuous bijection for which its inverse is also continuous.
\end{problem}

\begin{problem}
Let $B_{1}, B_{2}$ be two Borel sets of $\mathbb{R}$ and let $g:B_{1} \rightarrow B_{2}$ be a homeomorphism (that is, $g$ is a continuous bijection whose inverse is also continuous).  Show that if $B \subseteq B_{1}$ is any Borel set, then $g(B)$ is also a Borel set.
\end{problem}

\begin{problem}
Coming back to the function $f$ from Problem 8, show that $\lambda(f(C)) = 1$, where $C$ is the Cantor set.
\end{problem}
Hint:  What is $\lambda(f([0,1]\smallsetminus C))$?

\begin{problem}
Since $\lambda(f(C)) = 1 > 0$, explain why there exists $\Omega \subseteq f(C)$ such that $\Omega \notin \mathcal{L}(\mathbb{R})$.  Then, set $A = f^{-1}(\Omega) \subseteq [0,1]$.  Show that $A \in \mathcal{L}(\mathbb{R}) \smallsetminus \mathcal{B}(\mathbb{R})$.
\end{problem}



\end{document}









